\input{header}

\title{Peer Review \& Presentation Practice}
\subtitle{ECON 692 -- Applied Economics Seminar}
\author{Andrew Hobbs}
\institute{University of San Francisco}
\date{April 23, 2026}

\begin{document}

\begin{frame}[plain]
  \titlepage
\end{frame}

% ══════════════════════════════════════════════════════════
%  Agenda & Check-in
% ══════════════════════════════════════════════════════════

\begin{frame}{Today's Agenda}
  \small
  \begin{tabular}{@{}r@{\hspace{8pt}}l@{\hspace{16pt}}l@{}}
    \toprule
    \textbf{Time} & \textbf{Duration} & \textbf{Activity} \\
    \midrule
    4:35 & 10 min & Check-in \\
    4:45 & 10 min & How to give good peer review \\
    4:55 & 60 min & Peer review workshop (2 papers $\times$ 30 min) \\
    5:55 & 10 min & Break \\
    6:05 & 10 min & Presentation practice: setup \& norms \\
    6:15 & 50 min & Small-group presentation runs: round 1 \\
    7:05 & 50 min & Small-group presentation runs: round 2 \\
    7:55 & 10 min & Quick sprint: one thing to fix tonight \\
    8:05 & 10 min & Wrap-up \& next steps \\
    \bottomrule
  \end{tabular}
\end{frame}

\begin{frame}{Check-in}
  \textbf{\textcolor{usfgreen}{Full Draft is due tomorrow (Friday, April 24).}}

  \vspace{8pt}

  Quick show of hands:

  \vspace{4pt}

  \begin{itemize}
    \item Who has a complete draft to swap today?
    \item Who is still missing a section (intro, results, conclusion)?
    \item Who has slides started for final presentations?
  \end{itemize}

  \vspace{8pt}

  Today's goal: leave with a marked-up draft, a first dry run of your presentation, and a concrete list of what to fix before submission.
\end{frame}

% ══════════════════════════════════════════════════════════
\section{How to Give Good Peer Review}
% ══════════════════════════════════════════════════════════

\begin{frame}{What Peer Review Is For}
  \small
  The goal isn't to catch every typo. It's to tell the author what a \textbf{fresh reader} experiences.

  \vspace{6pt}

  \begin{columns}[T]
    \begin{column}{0.48\textwidth}
      \textbf{\textcolor{usfgreen}{Useful feedback}}
      \begin{itemize}\setlength\itemsep{1pt}
        \item ``I got lost in section 3 because\ldots''
        \item ``The key result wasn't clear until page 8.''
        \item ``I didn't understand why X identifies Y.''
        \item ``Figure 2 contradicts what the text says.''
      \end{itemize}
    \end{column}
    \begin{column}{0.48\textwidth}
      \textbf{\textcolor{red!60}{Less useful feedback}}
      \begin{itemize}\setlength\itemsep{1pt}
        \item ``Looks good.''
        \item ``Fix the typo on line 47.'' (without more)
        \item ``Add more citations.''
        \item ``Use different colors.''
      \end{itemize}
    \end{column}
  \end{columns}

  \vspace{6pt}

  The author already knows their paper. They need to know how it \emph{reads}.
\end{frame}

\begin{frame}{The Reviewer's Three Jobs}
  \begin{enumerate}
    \item \textbf{Summarize the paper back to the author.} In 2--3 sentences: what's the question, what's the method, what's the finding? If you can't, the author has a framing problem.

    \vspace{4pt}

    \item \textbf{Identify the two or three biggest issues.} Not twenty small ones. Major things first: is the question clear? Is the identification credible? Does the evidence match the claim?

    \vspace{4pt}

    \item \textbf{Suggest concrete fixes.} ``The intro buries the finding --- move paragraph 4 to the top'' is actionable. ``The intro is weak'' is not.
  \end{enumerate}

  \vspace{6pt}

  Do this first. Line edits come last, if at all.
\end{frame}

\begin{frame}{What to Look For: Checklist}
  \footnotesize
  \begin{columns}[T]
    \begin{column}{0.48\textwidth}
      \textbf{\textcolor{usfgreen}{Framing}}
      \begin{itemize}\setlength\itemsep{0pt}
        \item Is the question clear by the end of page 1?
        \item Why should I care?
        \item What's the contribution vs.\ existing work?
      \end{itemize}

      \vspace{4pt}

      \textbf{\textcolor{usfgreen}{Identification}}
      \begin{itemize}\setlength\itemsep{0pt}
        \item Is the empirical strategy stated clearly?
        \item What's the identifying assumption?
        \item What would make the skeptic doubt it?
      \end{itemize}
    \end{column}
    \begin{column}{0.48\textwidth}
      \textbf{\textcolor{usfgreen}{Evidence}}
      \begin{itemize}\setlength\itemsep{0pt}
        \item Do the figures/tables actually show what the text claims?
        \item Are units, samples, and sources labeled?
        \item Is there a ``money figure'' that sticks?
      \end{itemize}

      \vspace{4pt}

      \textbf{\textcolor{usfgreen}{Flow}}
      \begin{itemize}\setlength\itemsep{0pt}
        \item Where did you get lost?
        \item What paragraph felt like filler?
        \item What question did you want answered earlier?
      \end{itemize}
    \end{column}
  \end{columns}
\end{frame}

\begin{frame}{Ground Rules}
  \begin{enumerate}
    \item \textbf{Be direct, be kind.} Everyone here \emph{wants} the hard feedback --- that's why we're doing this a day before submission.

    \vspace{4pt}

    \item \textbf{Read like a reader, not a grader.} You're not assigning a score; you're telling the author how the paper reads.

    \vspace{4pt}

    \item \textbf{Write things down.} Verbal feedback evaporates. Margin notes + a short summary memo is what the author takes home.

    \vspace{4pt}

    \item \textbf{Don't defend your own paper today.} When you're the author, take notes and ask clarifying questions.
  \end{enumerate}
\end{frame}

% ══════════════════════════════════════════════════════════
\section{Final Presentation Date Assignments}
% ══════════════════════════════════════════════════════════

\begin{frame}{Final Presentation Date Assignments}
  \begin{columns}[T]
    \begin{column}{0.48\textwidth}
      \textbf{\textcolor{usfgreen}{May 7 (3 presenters)}}
      \begin{itemize}\setlength\itemsep{1pt}
        \item Tuan Truong
        \item Chris Tsang
        \item Evan Bruno
      \end{itemize}
    \end{column}
    \begin{column}{0.48\textwidth}
      \textbf{\textcolor{usfgreen}{May 14 (10 presenters)}}
      \footnotesize
      \begin{itemize}\setlength\itemsep{0pt}
        \item Khaliun Battogtokh
        \item Cecilia Vigil
        \item Tanisha Gauns
        \item Dia Talwar
        \item Jacob Guzman
        \item Austin Coffelt
        \item Ashley Thompson
        \item Amrit Gill
        \item Pawoumondom Alai
        \item Gaziz Makhanov
      \end{itemize}
    \end{column}
  \end{columns}
\end{frame}

% ══════════════════════════════════════════════════════════
\section{Peer Review Workshop}
% ══════════════════════════════════════════════════════════

\begin{frame}{Workshop Structure}
  \textbf{Two rounds --- 30 minutes each. You review \emph{two} other drafts; two people review yours.}

  \vspace{8pt}

  \begin{itemize}\setlength\itemsep{2pt}
    \item Each round: 25 min read + write memo, 5 min debrief with author
    \item I'll post pairings --- check the board before round 1 starts
    \item Swap drafts (paper or PDF); margin notes are fine
    \item Memo target: one short page (see next slide)
    \item Debrief: author takes notes, doesn't defend
  \end{itemize}
\end{frame}

\begin{frame}{The Review Memo}
  \small
  Hand the author a short memo (handwritten is fine) with four parts:

  \vspace{4pt}

  \begin{enumerate}\setlength\itemsep{2pt}
    \item \textbf{One-paragraph summary.} Question, method, finding, in your own words.

    \item \textbf{Two or three major comments.} The biggest things that would improve the paper. Be specific about where and what.

    \item \textbf{A handful of minor comments.} Unclear sentences, missing labels, figures that don't work, citations to add. Bullet list is fine.

    \item \textbf{One thing the paper does well.} Something to keep --- this helps the author know what not to cut.
  \end{enumerate}

  \vspace{6pt}

  Target: half a page to one page, $\sim$5 minutes to write after reading.
\end{frame}

\begin{frame}{Reading Strategy (25 min per paper is tight)}
  \footnotesize
  You've got 25 minutes to read and write before the 5-minute debrief. Don't try to read every word.

  \vspace{6pt}

  \begin{enumerate}\setlength\itemsep{2pt}
    \item \textbf{First 3 min: skim.} Read the intro, the headings, look at the figures. Write down what you think the paper is about \emph{before} reading it carefully.

    \item \textbf{Next 12 min: read.} Focus on the intro, empirical strategy, and results sections. Skim the rest. Mark confusing passages as you go.

    \item \textbf{Next 5 min: look at evidence.} Does each figure/table support the claim in the text? Are the units, samples, and sources clear?

    \item \textbf{Last 5 min: write the memo.} Summary first, then major comments, then minor comments, then the thing it does well.

    \item \textbf{Debrief (5 min, separate block).} Walk the author through your memo. They take notes; they don't defend.
  \end{enumerate}

  \vspace{4pt}

  Prioritize the parts that drive the argument.
\end{frame}

% ══════════════════════════════════════════════════════════
\section{Presentation Practice}
% ══════════════════════════════════════════════════════════

\begin{frame}{Presentation Reminder: The Arc}
  \footnotesize
  15 minutes of talking $+$ Q\&A. Target $\sim$12--15 slides.

  \vspace{4pt}

  \begin{columns}[T]
    \begin{column}{0.48\textwidth}
      \begin{enumerate}\setlength\itemsep{0pt}
        \item Title
        \item Motivation (the hook)
        \item Research question
        \item Related work (1 slide)
        \item Data
        \item Summary stats
        \item Empirical strategy
      \end{enumerate}
    \end{column}
    \begin{column}{0.48\textwidth}
      \begin{enumerate}\setlength\itemsep{0pt}
        \setcounter{enumi}{7}
        \item Identification
        \item Main result
        \item Robustness / mechanisms
        \item Interpretation
        \item Conclusion
      \end{enumerate}

      \vspace{4pt}

      Spend most of your time on 7--10. Don't burn 8 minutes on background.
    \end{column}
  \end{columns}
\end{frame}

\begin{frame}{What a Good Practice Run Looks Like}
  \begin{enumerate}
    \item \textbf{Talk through it end to end.} No stopping to fix slides. Finish, then revise. Talking while you edit is a trap.

    \vspace{4pt}

    \item \textbf{Time yourself.} If you're at 20 minutes, cut. If you're at 8, add. Don't guess.

    \vspace{4pt}

    \item \textbf{Practice the transitions.} The hardest part of a talk isn't any single slide --- it's the sentence that moves you from one to the next.

    \vspace{4pt}

    \item \textbf{Practice the first 60 seconds out loud.} The opening is the most-rehearsed part of a good talk. Know the first two sentences cold.
  \end{enumerate}
\end{frame}

\begin{frame}{Giving Feedback on a Practice Talk}
  \small
  When you're the audience, take notes on:

  \vspace{4pt}

  \begin{columns}[T]
    \begin{column}{0.48\textwidth}
      \textbf{\textcolor{usfgreen}{Content}}
      \begin{itemize}\setlength\itemsep{0pt}
        \item Did the question land?
        \item Did the identification story hold?
        \item Did you believe the main result?
        \item What question do you want to ask?
      \end{itemize}
    \end{column}
    \begin{column}{0.48\textwidth}
      \textbf{\textcolor{usfgreen}{Delivery}}
      \begin{itemize}\setlength\itemsep{0pt}
        \item Pace: too fast, too slow?
        \item Was the main figure legible?
        \item Were slide titles claims or labels?
        \item Filler words? (``um,'' ``like,'' ``basically'')
      \end{itemize}
    \end{column}
  \end{columns}

  \vspace{6pt}

  Give the presenter \textbf{two things to keep} and \textbf{two things to change} before the real talk.
\end{frame}

\begin{frame}{Small-Group Presentation Runs}
  \textbf{Pairs. Two rounds with different partners.} Each person gets 20 minutes per round: 15 min talk + 5 min feedback.

  \vspace{8pt}

  \textbf{Ground rules:}
  \begin{itemize}\setlength\itemsep{1pt}
    \item \textbf{Assign a timekeeper} in each pair --- they call time at 15 min
    \item Use whatever slides you have --- even if it's just an outline
    \item Stand up and talk through it. Don't read off the screen.
    \item The listener takes notes
    \item No interruptions during the talk; all feedback at the end
  \end{itemize}

  \vspace{8pt}

  If your slides aren't built yet, practice with a whiteboard or just talk through the structure. The point is to hear yourself say it out loud.
\end{frame}

\begin{frame}{Round 1 vs.\ Round 2}
  \begin{columns}[T]
    \begin{column}{0.48\textwidth}
      \textbf{\textcolor{usfgreen}{Round 1 (6:15--7:05)}}
      \begin{itemize}\setlength\itemsep{1pt}
        \item First partner assignment
        \item Full 15-minute talk
        \item Listener: focus on \textbf{content} --- did the question and result land?
      \end{itemize}
    \end{column}
    \begin{column}{0.48\textwidth}
      \textbf{\textcolor{usfgreen}{Round 2 (7:05--7:55)}}
      \begin{itemize}\setlength\itemsep{1pt}
        \item New partner (I'll reshuffle)
        \item Incorporate what you learned in round 1
        \item Listener: focus on \textbf{delivery} --- pacing, transitions, legibility
      \end{itemize}
    \end{column}
  \end{columns}

  \vspace{10pt}

  Round 2 is where the real gains are. The second time you say something out loud is always better than the first.
\end{frame}

% ══════════════════════════════════════════════════════════
\section{Work Sprint}
% ══════════════════════════════════════════════════════════

\begin{frame}{Quick Sprint: One Thing to Fix Tonight (7:55--8:05)}
  10 minutes. Pick the \textbf{single most useful thing} you can fix tonight based on what you just heard.

  \vspace{8pt}

  \begin{columns}[T]
    \begin{column}{0.48\textwidth}
      \textbf{\textcolor{usfgreen}{Draft}}
      \begin{itemize}\setlength\itemsep{1pt}
        \item The biggest framing issue a reviewer raised
        \item A confusing figure caption
      \end{itemize}
    \end{column}
    \begin{column}{0.48\textwidth}
      \textbf{\textcolor{usfgreen}{Slides}}
      \begin{itemize}\setlength\itemsep{1pt}
        \item Cut the slide everyone fumbled on
        \item Rewrite your title + question slide
      \end{itemize}
    \end{column}
  \end{columns}

  \vspace{10pt}

  Write the one thing down. Don't start it here --- just commit to finishing it tonight.
\end{frame}

% ══════════════════════════════════════════════════════════
\section{Wrap-up}
% ══════════════════════════════════════════════════════════

\begin{frame}{Next Week \& Deadlines}
  \small
  \textbf{This week:}
  \begin{itemize}\setlength\itemsep{0pt}
    \item Weekly progress report due \textbf{Wednesday}
    \item \textbf{\textcolor{usfgreen}{Full Draft due tomorrow (April 24)}}
  \end{itemize}

  \textbf{Next week (4/30):} Revision and polishing. Come in with reviewer comments from today and a plan for what to tackle.

  \textbf{Presentations:}
  \begin{itemize}\setlength\itemsep{0pt}
    \item \textbf{May 7:} first group ($\sim$3 presenters)
    \item \textbf{May 14:} remaining presenters + reproducibility package due 5/15
  \end{itemize}

  \begin{block}{Before You Leave Tonight}
    Submit tomorrow. Even if it's imperfect --- a complete draft you can revise beats a perfect intro with nothing behind it.
  \end{block}
\end{frame}

\end{document}
